\section{Uniform cancellation with a surviving source}\label{sec:mixed}

We prove Theorems~\ref{thm:intromixed} and~\ref{thm:relaxation}. The extra source cohomology
prevents direct use of \eqref{eq:Xcontract}. The block estimate of
Section~\ref{sec:polar} controls the resulting term without an inverse
of the surviving orbit operator.

\subsection{Equations in an adapted orthonormal splitting}
For any surjective $D$, choose orthonormal frames adapted to
$V=\ker D\oplus(\ker D)^\perp$, and write
\[
 D=[0,F],\qquad A_j=[H_j,T_j],\qquad
 \Phi_V=\begin{pmatrix}\Phi_{HH}&\Phi_{HA}\\
                         \Phi_{AH}&\Phi_{AA}\end{pmatrix}.
\]
Add to the source frame velocity the skew-Hermitian generator
\[
 K_V=\frac12\begin{pmatrix}0&\Phi_{HA}\\-\Phi_{AH}&0\end{pmatrix}.
\]
The off-diagonal blocks have sizes $n_s$ by $r$ and $r$ by $n_s$;
no equality of their sizes is needed. The kernel component of $\dot D$
remains zero: its unrotated term
$F\Phi_{AH}/2$ is canceled by $F(K_V)_{AH}$.
Equation~\eqref{eq:mapflow} then becomes
\begin{align}
 \dot F&=\tfrac12(F\Phi_{AA}-\Phi_UF),&
 \dot H_j&=\tfrac12(H_j\Phi_{HH}-\Phi_WH_j),\label{eq:adapted1}\\
 \dot T_j&=H_j\Phi_{HA}
              +\tfrac12(T_j\Phi_{AA}-\Phi_WT_j).\label{eq:adapted2}
\end{align}
In particular, with $X_j=T_jF^{-1}$,
\begin{equation}\label{eq:mixedX}
 \dot X_j=\tfrac12(X_j\Phi_U-\Phi_WX_j)+H_j\Phi_{HA}F^{-1}.
\end{equation}
The second term is the coupling absent from the invertible-source
argument.

For this subsection put $B=F^\dagger F$, $L=FF^\dagger$,
$C_T=\sum_jT_j^\dagger T_j$,
$C_s=\sum_jH_j^\dagger H_j$ and $C_w=\sum_jH_jH_j^\dagger$.
The exact coefficients are
\begin{equation}\label{eq:adaptedcoeff}
 \begin{split}
 \Omega_V&=
 \begin{pmatrix}
 z_0\id+C_s&E_R\\
 E_R^\dagger&z_0e^{-B}+K_B(C_T)
 \end{pmatrix},
 \qquad E_R=\sum_jH_j^\dagger T_jR(B),\\
 \Omega_U&=J_L(FC_TF^\dagger)-z_0e^{-L},\\
 \Omega_W&=z_1\id-C_w-Y,\qquad
 Y=\sum_jX_j(\id-e^{-L})X_j^\dagger\geq0 .
 \end{split}
\end{equation}
Define
\begin{equation}\label{eq:qslambda}
 q=\sum_j\norm{H_j}_{\HS}^2,\qquad
 s=\sum_j\norm{X_j}_{\HS}^2,\qquad
 \lambda=\lambda_{\min}(F^\dagger F).
\end{equation}
Since $\norm{FR(B)}\leq\lambda^{-1/2}$,
Cauchy--Schwarz gives
\begin{equation}\label{eq:mixedblockbound}
 \norm{E_R}_{\HS}\leq\sqrt{qs/\lambda}.
\end{equation}
Also $\norm{Y}_{\HS}\leq\tr Y\leq s$.
The source compression $\Phi_{AA}$ is positive and $\Phi_U$ is
negative, so the least-singular-value calculation
\eqref{eq:lambdaderivative}, now applied to $F$, shows that
$\lambda$ is nondecreasing.

\subsection{Uniform bounds for the surviving arrows}
Allow arbitrary surviving tuples and charges in the strict upper
half-plane. Put $b_i=\Ima z_i>0$, $a=\arg z_1\in(0,\pi)$ and define
the ordinary cohomology phases
\[
 \phi_s=\arg(z_0\id+C_s),\qquad
 \phi_w=\arg(z_1\id-C_w),\qquad
 \delta_s=\Phi_{HH}-\phi_s,\quad \delta_w=\Phi_W-\phi_w.
\]
These ordinary coefficients are normal. Their source singular values
are at least $b_0$. Lemma~\ref{lem:partialpolar} gives
\begin{equation}\label{eq:delta}
 \norm{\delta_s}_{\HS}\leq\frac{4\pi}{b_0}\sqrt{qs/\lambda},\qquad
 \norm{\delta_w}_{\HS}\leq s/b_1
\end{equation}
when $\sqrt{qs/\lambda}/b_0\leq1/8$.
For the target inequality, the scalar function
$g(x)=\arg(z_1-x)$ has derivative at most $1/b_1$.
Expanding two Hermitian matrices in their separate eigenbases shows
$\norm{g(C_w+Y)-g(C_w)}_{\HS}\leq\norm Y_{\HS}/b_1$.
No operator monotonicity assertion is needed.

The unconditional off-diagonal estimate
$\norm{\Phi_{HA}}_{\HS}\leq(2\pi/b_0)\sqrt{qs/\lambda}$,
together with $\Phi_W\geq a\id$, gives from \eqref{eq:mixedX}
\begin{equation}\label{eq:sineq}
 s'\leq\left(-a+\frac{4\pi q}{b_0\lambda}\right)s.
\end{equation}
To bound $q$, use the ordinary cohomology mass
\begin{equation}\label{eq:cohomass}
 \mathcal M_H=\norm{z_0\id+C_s}_1+\norm{z_1\id-C_w}_1.
\end{equation}
The scalar inequalities $|z_0+x|\geq x+\Rea z_0$ and
$|z_1-x|\geq x-\Rea z_1$, together with
$\tr C_s=\tr C_w=q$, imply
\begin{equation}\label{eq:qfrommass}
 q\leq\tfrac12(\mathcal M_H-n_s\Rea z_0+n_w\Rea z_1).
\end{equation}
For the ordinary degree-zero phase equation, its derivative in source
and target phase eigenbases is
\[
 \sum_{j,i,l}|(H_j)_{li}|^2
       (\cos\alpha_i-\cos\beta_l)(\alpha_i-\beta_l)\leq0,
\]
because cosine decreases on $(0,\pi)$.
Equivalently, its arrow gradient is
$2(H_j\cos\phi_s-\cos\phi_w H_j)$, with total
Hilbert--Schmidt norm at most $4\sqrt q$.
The additional arrow velocity in \eqref{eq:adapted1} is
$(H_j\delta_s-\delta_wH_j)/2$, so
\begin{equation}\label{eq:cohomassineq}
 \mathcal M_H'\leq2q(\norm{\delta_s}_{\HS}+\norm{\delta_w}_{\HS}).
\end{equation}
Only this auxiliary ordinary mass, not the full complex mass, has
been used in a dissipation inequality.

Let $M_0=\mathcal M_H(0)$ and put
\begin{equation}\label{eq:Qgeneral}
 Q=\tfrac12(M_0-n_s\Rea z_0+n_w\Rea z_1)+1.
\end{equation}
Then $Q\geq q(0)+1\geq1$.

\begin{proposition}\label{prop:bootstrap}
Suppose $\lambda_0=\sigma_{\min}(F_0)^2$ satisfies
\begin{equation}\label{eq:threshold}
 \lambda_0\geq L_*:=\max\left\{
 1,\frac{8\pi Q}{ab_0},\frac{8\sqrt Q M}{b_0},
 \frac{64\pi Q^{3/2}M/b_0+8QM^2/b_1}{a}\right\}.
\end{equation}
For all $t\geq0$, one has $q(t)\leq Q-3/4$ and
\begin{align}
 s(t)&\leq\frac{M^2}{\lambda_0}e^{-at/2},\label{eq:smixed}\\
 \norm{\delta_s}_{\HS}
 &\leq\frac{4\pi\sqrt Q M}{b_0\lambda_0}e^{-at/4},
 &\norm{\delta_w}_{\HS}
 &\leq\frac{M^2}{b_1\lambda_0}e^{-at/2}.\label{eq:deltamixed}
\end{align}
\end{proposition}

\begin{proof}
Initially $q\leq Q-1$ and $s\leq M^2/\lambda_0$, since
$\norm{F_0^{-1}}=\lambda_0^{-1/2}$.
While $q\leq Q$, monotonicity of $\lambda$ and
\eqref{eq:threshold} give $s'\leq-as/2$.
This proves \eqref{eq:smixed} on that interval and ensures
$\sqrt{qs/\lambda}/b_0\leq1/8$.
Equations~\eqref{eq:delta} give \eqref{eq:deltamixed}.
Integrating \eqref{eq:cohomassineq} yields
\[
 \mathcal M_H(t)\leq M_0+
 \frac{32\pi Q^{3/2}M/b_0+4QM^2/b_1}{a\lambda_0}
 \leq M_0+\tfrac12 .
\]
Equation~\eqref{eq:qfrommass} gives $q\leq Q-3/4<Q$, closing
the bootstrap by continuity. Global existence extends the bounds to
all $t\geq0$.
\end{proof}

Take $\Lambda_*=\sqrt{L_*}$. If $q(0)\leq R$, the triangle inequality
gives $M_0\leq2R+n_s|z_0|+n_w|z_1|$. The same proof therefore works
with $Q$ replaced by the uniform upper bound
\begin{equation}\label{eq:QR}
 Q_R=R+\tfrac12\bigl[n_s(|z_0|-\Rea z_0)
                         +n_w(|z_1|+\Rea z_1)\bigr]+1.
\end{equation}
This gives a threshold uniform on initial surviving-arrow norm balls,
independent of acyclic rank, arrow count and the largest singular value
of $F_0$. Charge dependence is explicit in \eqref{eq:threshold}; no
uniformity near the boundary of the upper half-plane is asserted.

\subsection{Growth and cancellation of the acyclic sector}
Put $\sigma=a/4$ and $\kappa=\pi/(2e)$.
For the trace-norm constants put
\begin{equation}\label{eq:mixedsourceconstant}
 d=\min(n_s,r),\qquad A_0=2\sqrt{dQ}M+r\kappa M^2.
\end{equation}
Compare the source coefficient in \eqref{eq:adaptedcoeff} with
\[
 \Omega_{V,0}=\diag(z_0\id+C_s,z_0e^{-B}).
\]
Its error has trace norm at most
\begin{equation}\label{eq:sourceerror}
 \norm{\Omega_V-\Omega_{V,0}}_1
 \leq\frac{A_0}{\lambda_0}e^{-\sigma t}.
\end{equation}
Indeed the two off-diagonal blocks contribute at most
$2\sqrt d\norm{E_R}_{\HS}$, since each has rank at most $d$.
Lemma~\ref{lem:heatbounds}, with rank $r$, applies to $T_j=X_jF$
on the acyclic block and gives
$\norm{K_B(C_T)}_1\leq r\kappa s$.
Similarly the degree-one error relative to $-z_0e^{-L}$ is at most
$rM^2\lambda_0^{-1}e^{-at/2}$.

Apply Lemma~\ref{lem:weightedpolar} to a vector in the acyclic source
block realizing $\lambda$, and to its image under $F/\sqrt\lambda$.
The reference singular value is $|z_0|e^{-\lambda}$ in both cases.
The reference phase on that source direction is $\arg z_0$,
regardless of the phase on the surviving source block.
Consequently
\begin{equation}\label{eq:mixedgapineq}
 \lambda'\geq\lambda
       \left[\pi-C_1e^{(\lambda-\sigma t)/2}\right],
\end{equation}
where the following constant has no condition-number dependence:
\[
 C_1=\frac{\pi}{\sqrt{|z_0|}}
 \left(\sqrt{A_0}+\sqrt r M\right).
\]
We used $\lambda_0\geq1$ to discard a favorable factor.
The same estimate holds for every vector in a repeated least
eigenspace, so it also controls the one-sided derivatives of $\lambda$.
Set $L_0=\max\{0,2\log(2C_1/\pi)\}$, with $L_0=0$ if $C_1=0$.
Because $\lambda_0\geq1$, the line $\sigma t-L_0$ lies below
$\lambda$ up to time $(1+L_0)/\sigma$. At a later contact point,
\eqref{eq:mixedgapineq} gives
$\lambda'\geq\pi\lambda/2\geq\pi/2>2\sigma=a/2$.
The same barrier comparison as before gives
\begin{equation}\label{eq:mixedgap}
 \lambda(t)\geq\max\{\lambda_0,\sigma t-L_0\}.
\end{equation}
In particular,
\begin{equation}\label{eq:mixedheat}
 \norm{e^{-B}}=\norm{e^{-L}}
 \leq\min(e^{-\lambda_0},e^{L_0}e^{-\sigma t})
 \leq\frac{2e^{L_0/2}}{e\lambda_0}e^{-\sigma t/2}.
\end{equation}
For the last step, take the geometric mean and use
$\lambda_0e^{-\lambda_0/2}\leq2/e$.

Define the instantaneous cohomology coefficient to have blocks
$z_0\id+C_s$ and $z_1\id-C_w$. Equations
\eqref{eq:adaptedcoeff}, \eqref{eq:smixed},
\eqref{eq:sourceerror} and \eqref{eq:mixedheat} show that
\begin{equation}\label{eq:mixedcancelconstant}
 C_{\mathrm{can}}=
 2\left[A_0+(r+1)M^2+\frac{4r|z_0|}{e}e^{L_0/2}\right]
\end{equation}
satisfies
\[
 \norm{\Omega-(\Omega_H\oplus0_{\ac})}_1
 \leq\tfrac12 C_{\mathrm{can}}\lambda_0^{-1}e^{-at/8}.
\]
Here the two acyclic heat blocks each have rank $r$, and the target
error has trace norm at most $s$. These account for every term in
\eqref{eq:mixedcancelconstant}.
The trace norm is $1$-Lipschitz with respect to itself, so the same
bound holds for $|m-\mathcal M_H|$.
Adding the bounds and using $\lambda_0\geq\Lambda^2$ proves the
quantitative cancellation assertion in Theorem~\ref{thm:intromixed}.
Taking $C=\max\{C_{\mathrm{can}},L_0\}$ supplies one constant for its
growth and cancellation conclusions. All estimates use tuple norms and
only the least acyclic singular value; no stability hypothesis has entered.

\subsection{Relaxation from finite initial energy}
For the rest of this section assume $n_s=n_w=n$ and
$z_0=-1+\ii$, $z_1=1+\ii$, so $a=\pi/4$ and $b_0=b_1=1$.
The surviving tuple need not be initially balanced.
In fixed holomorphic coordinates, $\ker D$ is fixed and
$A_j|_{\ker D}=B_j$. Its induced metrics are the restriction of
$h_V$ and the metric $h_W$. The restriction's metric velocity is
$-\Phi_{HH}$, as represented by \eqref{eq:adapted1}.

Multiply both induced metrics by $e^{\pi t/2}$ and call them
$h_s,h_w$. This leaves the arrow norms unchanged. The common scalar
contributions to the logarithmic determinants cancel because the
surviving dimensions agree. Using the initial orthonormal determinant
volume forms, define on all positive metric pairs
\begin{equation}\label{eq:KN}
 \mathcal K=q+\log\det h_s-\log\det h_w,\qquad
 \mu_s=C_s-\id,\quad\mu_w=C_w-\id.
\end{equation}
Assume $\mathcal K_*:=\inf\mathcal K> -\infty$ and put
$E_0=\mathcal K(0)-\mathcal K_*$. The infimum need not be attained.

The normalized metric velocities are
\[
 \xi_s=\arctan\mu_s-\delta_s,\qquad
 \xi_w=-\arctan\mu_w-\delta_w.
\]
Variation of \eqref{eq:KN} gives
\[
 \mathcal K'
 =-\tr(\mu_s\arctan\mu_s+\mu_w\arctan\mu_w)
      +\tr(\mu_s\delta_s-\mu_w\delta_w).
\]
Here $M_0\geq2n$, so $Q\geq2n+1\geq3$. Since $q\leq Q$,
the moment spectra lie in $[-1,Q-1]$. Set $J_Q=1+(Q-1)^2$.
Integrating $(\arctan x)'=(1+x^2)^{-1}$ from zero gives
$x\arctan x\geq x^2/J_Q$ on this interval.
With squared norms summing source and target Hilbert--Schmidt squares,
Young's inequality yields
\begin{equation}\label{eq:KNderivative}
 \mathcal K'\leq-\frac1{2J_Q}\norm\mu^2+\frac{J_Q}{2}\norm\delta^2.
\end{equation}
By \eqref{eq:deltamixed},
\[
 \int_0^\infty\norm{\delta(t)}^2\,dt
 \leq\frac{32\pi^2QM^2+M^4}{a\lambda_0^2}.
\]
Consequently, with
\begin{equation}\label{eq:mixedenergyconstant}
 C_{\mathrm{en}}(Q,M)=\frac{J_Q}{2a}(32\pi^2QM^2+M^4),
\end{equation}
one has
\begin{equation}\label{eq:KNbound}
 \begin{gathered}
 0\leq\mathcal K(t)-\mathcal K_*,\\
 \mathcal K(t)-\mathcal K_*+
 \frac1{2J_Q}\int_0^t\norm{\mu(u)}^2\,du
 \leq E_0+C_{\mathrm{en}}(Q,M)\lambda_0^{-2}
 \leq E_0+C_{\mathrm{en}}(Q,M)\Lambda^{-4}.
 \end{gathered}
\end{equation}
Only the lower bound for $\mathcal K$ is used. A bound on $E_0$ is
needed for a uniform numerical energy bound across initial metrics;
its finiteness suffices for the following qualitative conclusion.

The adapted equations and phase bounds give
\[
 \sum_j\norm{\dot H_j}_{\HS}^2\leq\pi^2q,\qquad
 \norm{\dot\mu_s}_{\HS},\ \norm{\dot\mu_w}_{\HS}\leq2\pi q.
\]
Thus $\mu$ is uniformly continuous, independently of $k$.
An integrable nonnegative squared norm with this property tends to zero:
otherwise disjoint intervals of a fixed positive length would each
contribute a fixed positive amount to its integral.
It follows that $\mu_s,\mu_w\to0$, hence $\mathcal M_H\to2n$.
The cancellation estimate then gives $m(h(t))\to2n$.
This proves Theorem~\ref{thm:relaxation}.
When $M=0$, the forcing $\delta$ vanishes, but $\mu$ need not vanish
at finite times for nonbalanced initial metrics.

\subsection{Balanced minima and an unattained infimum}
Suppose $(h_s^*,h_w^*)$ is a balanced metric pair.
In orthonormal frames for that pair, a geodesic of positive metrics
has the form $(e^{\tau S_s},e^{\tau S_w})$.
After diagonalizing $S_s,S_w$, the squared arrow norm is a sum of
nonnegative multiples of $e^{\tau(\lambda_w-\lambda_s)}$;
the logarithmic terms are affine. Thus $\mathcal K$ is convex along
the geodesic, and balance makes its derivative zero at $\tau=0$.
Every positive pair is the other endpoint of such a geodesic, so
this balanced pair attains the global minimum
\[
 \mathcal K_*=n+\log\det h_s^*-\log\det h_w^*.
\]
In particular, initially balanced orthonormal metrics give
$M_0=2n$, $Q=2n+1$, $J_Q=1+4n^2$ and $E_0=0$.
Equations \eqref{eq:threshold} and \eqref{eq:KNbound} recover
Corollary~\ref{cor:balanced}, with threshold depending only on $n,M$.
No strict convexity or positive stabilizer gap is required.

The finite-infimum hypothesis also allows nonattainment.
Take $B_1=\id_2$ and $B_2=e_{12}$, the matrix with a single nonzero
entry $1$ in position $(1,2)$.
For arbitrary positive metrics put
$X=h_s^{-1/2}h_wh_s^{-1/2}$. Then
\[
 \mathcal K=\tr X-\log\det X+\norm{B_2}_{\HS,h}^2>2,
\]
since $x-\log x\geq1$ and $B_2\ne0$.
Along $h_s=h_w=\diag(e^{-t},e^t)$, however,
$\mathcal K=2+e^{-2t}$. Its infimum is $2$ and is not attained.
Theorem~\ref{thm:relaxation} applies to this surviving representation
when the initial acyclic scale is sufficiently large.

Cancellation alone cannot imply relaxation to balance: for zero
surviving arrows with these charges, $\mathcal M_H=2\sqrt2\,n$ and
$\mu_s=\mu_w=-\id$, while $\mathcal K$ has no finite lower bound.
The exponential comparison with instantaneous cohomology remains valid.
Even under the finite-infimum hypothesis, \eqref{eq:KNbound} supplies
neither a uniform final relaxation rate nor convergence of normalized
surviving metrics. It is not an inverse estimate for the Hessian of
$\mathcal K$ at a changing stabilizer.
